% ============================================================
%  Преамбула учебного пособия
%  Оформление по ГОСТ 7.60-2003, ГОСТ Р 7.0.100-2018
%  Компиляция: xelatex
% ============================================================

% ---------- Шрифты и языки ----------
\usepackage{fontspec}
\usepackage{polyglossia}
\setmainlanguage{russian}
\setotherlanguage{english}
\setmainfont{Times New Roman}
\setsansfont{Arial}
\setmonofont{Consolas}[Scale=0.92]

% ---------- Геометрия страницы (поля 2 см) ----------
\usepackage{geometry}
\geometry{a4paper,left=2cm,right=2cm,top=2cm,bottom=2cm}

% ---------- Интерлиньяж 1,5 и абзацный отступ 1,25 см ----------
\usepackage{setspace}
\onehalfspacing
\usepackage{indentfirst}
\setlength{\parindent}{1.25cm}

% ---------- Микротипографика: убирает вылеты за поля ----------
\usepackage{microtype}
% Длинные имена классов в моноширинном начертании (например
% UnsupportedOperationException или java.util.concurrent.atomic) не переносятся
% и выталкивают строку за поле. \sloppy разрешает TeX растягивать межсловные
% пробелы вместо этого: изредка строка выйдет разреженной, зато текст не полезет
% на поля. Для технической книги с обилием идентификаторов это верный размен.
\sloppy
\emergencystretch=4em
\hbadness=10000
\vbadness=10000

% Разрешаем разрывать длинные идентификаторы после точки и подчёркивания —
% без дефиса, чтобы разрыв не читался как часть имени.
\usepackage{xurl}

% ---------- Цвет, графика, таблицы ----------
\usepackage{xcolor}
\usepackage{graphicx}
\usepackage{booktabs}
\usepackage{longtable}
\usepackage{tabularx}
\usepackage{array}
\usepackage{multirow}
\newcolumntype{L}[1]{>{\raggedright\arraybackslash}p{#1}}
\newcolumntype{C}[1]{>{\centering\arraybackslash}p{#1}}

% ---------- Подписи по ГОСТ: «Таблица 1.1 — Название» ----------
\usepackage{caption}
% ГОСТ требует тире между номером и названием: «Таблица 1.1 — Название»
\DeclareCaptionLabelSeparator{emdash}{~--- }
\captionsetup{labelsep=emdash,font=normalsize,justification=centering}
\captionsetup[table]{position=above,justification=raggedright,singlelinecheck=false}
\captionsetup[figure]{position=below}
% подписи листингов настраиваются ниже, после загрузки minted

% ---------- Листинги кода ----------
%
%  Используется minted, а НЕ listings. Причина принципиальная: под XeLaTeX
%  пакет listings неверно разбирает кириллицу внутри кода — ASCII-знак
%  вплотную перед русской буквой переставляется в конец слова
%  («// Java-объектом» печатается как «// Javaобъектом-»), а строка,
%  начинающаяся с кириллицы, приклеивается к предыдущей и сбивает нумерацию.
%  Ни literate-подстановка, ни объявление байтов 0x80-0xFF буквами не лечат
%  это под XeLaTeX. minted передаёт код в Pygments, который корректно работает
%  с UTF-8. Проверено на test_minted.tex.
%
%  ВАЖНО: minted требует запуска xelatex с ключом -shell-escape и наличия
%  Python-пакетов pygments и latexminted (см. README).
%
\usepackage[newfloat]{minted}
\usepackage{float}

\definecolor{codebg}{RGB}{248,248,248}

\setminted{
  fontsize=\footnotesize,
  linenos,
  numbersep=8pt,
  frame=single,
  framesep=4pt,
  breaklines=true,
  breakanywhere=false,
  tabsize=4,
  bgcolor=codebg,
  autogobble=false
}

% Плавающее окружение листинга: подпись снизу, название по-русски
\SetupFloatingEnvironment{listing}{name=Листинг}
\SetupFloatingEnvironment{listing}{listname=Список листингов}
\captionsetup[listing]{position=below,justification=raggedright,
                       singlelinecheck=false}

% Листинг не должен начинаться в самом низу полосы: если до конца страницы
% осталось меньше четырёх строк, листинг целиком переносится на следующую.
\usepackage{needspace}

% ---------- Гиперссылки (в печати не мешают, в PDF работают) ----------
\usepackage[hidelinks,unicode]{hyperref}
\hypersetup{
  pdftitle={Современные технологии программирования},
  pdflang={ru-RU},
  bookmarksnumbered=true
}

% ---------- Заголовки глав и разделов ----------
\usepackage{titlesec}
\titleformat{\chapter}[display]
  {\normalfont\huge\bfseries\centering}
  {\chaptertitlename\ \thechapter}{14pt}{\Large}
\titlespacing*{\chapter}{0pt}{-20pt}{28pt}
\titleformat{\section}{\normalfont\Large\bfseries}{\thesection.}{0.6em}{}
\titleformat{\subsection}{\normalfont\large\bfseries}{\thesubsection.}{0.6em}{}
\titleformat{\subsubsection}{\normalfont\normalsize\bfseries}{\thesubsubsection.}{0.6em}{}

% ---------- Колонтитулы ----------
\usepackage{fancyhdr}
\pagestyle{fancy}
\fancyhf{}
\fancyhead[C]{\small\leftmark}
\fancyfoot[C]{\thepage}
\renewcommand{\headrulewidth}{0.4pt}
\fancypagestyle{plain}{\fancyhf{}\fancyfoot[C]{\thepage}\renewcommand{\headrulewidth}{0pt}}
\renewcommand{\chaptermark}[1]{\markboth{\thechapter.\ #1}{}}

% ---------- Списки поплотнее ----------
\usepackage{enumitem}
\setlist{topsep=4pt,itemsep=2pt,parsep=0pt}

% ---------- Врезки: «Важно», «Частая ошибка», «Аналогия» ----------
\usepackage[most]{tcolorbox}
\newtcolorbox{vrezka}[1][]{
  colback=blue!3!white,
  colframe=blue!45!black,
  boxrule=0.6pt,
  left=6pt,right=6pt,top=5pt,bottom=5pt,
  breakable,
  #1
}
\newtcolorbox{oshibka}[1][]{
  colback=red!3!white,
  colframe=red!50!black,
  boxrule=0.6pt,
  left=6pt,right=6pt,top=5pt,bottom=5pt,
  breakable,
  #1
}

% ---------- Окружения методического аппарата ----------
% Контрольные вопросы, тестовые задания, практикум
\newcommand{\praktikum}[1]{%
  \section*{Практикум}%
  \addcontentsline{toc}{section}{Практикум}%
  \markright{Практикум}%
  #1%
}
\newcommand{\kontrolvoprosy}{%
  \section*{Контрольные вопросы}%
  \addcontentsline{toc}{section}{Контрольные вопросы}%
}
\newcommand{\testzadaniya}{%
  \section*{Тестовые задания}%
  \addcontentsline{toc}{section}{Тестовые задания}%
}

% Тестовый вопрос: \vopros{стем}{вариант а}{вариант б}{вариант в}{вариант г}
% Используется внутри enumerate после \testzadaniya.
% Метки а)-г) проставляются явно: автоматических кириллических меток для
% второго уровня списка эта сборка не даёт (\asbuk здесь недоступен).
% samepage не позволяет странице разорвать вопрос между стемом и вариантами;
% \item намеренно оставлен снаружи samepage, чтобы разрывы между вопросами
% остались разрешёнными.
\newcommand{\vopros}[5]{%
  \item #1
  \begin{samepage}
  \begin{enumerate}
    \item[а)] #2
    \item[б)] #3
    \item[в)] #4
    \item[г)] #5
  \end{enumerate}
  \end{samepage}%
}

% ---------- Глубина оглавления и нумерации ----------
\setcounter{tocdepth}{2}
\setcounter{secnumdepth}{3}

% ---------- Нумерация рисунков/таблиц/листингов по главам ----------
\usepackage{chngcntr}
\counterwithin{figure}{chapter}
\counterwithin{table}{chapter}
% счётчик listing создаётся окружением minted позже, поэтому привязываем его
% к главам уже после загрузки всех пакетов
\AtBeginDocument{\counterwithin{listing}{chapter}}

% ---------- Термины глоссария (простое выделение) ----------
\newcommand{\term}[1]{\textbf{#1}}
\newcommand{\eng}[1]{\textit{\foreignlanguage{english}{#1}}}

% ============================================================
%  Диаграммы: рисуем средствами LaTeX, а не картинками.
%  Это даёт векторную графику, единый шрифт с текстом книги
%  и компактную вёрстку, влезающую в полосу набора.
% ============================================================
\usepackage{tikz}
\usepackage[edges]{forest}
\usepackage{pgfplots}
\pgfplotsset{compat=1.18}
\usetikzlibrary{arrows.meta,positioning,shapes.geometric,fit,backgrounds,calc,shapes.multipart}

% Дерево иерархии классов: растёт слева направо, компактное
\forestset{
  hierarchy/.style={
    for tree={
      grow'=0,
      draw,
      rounded corners=2pt,
      font=\ttfamily\footnotesize,
      edge={-Latex, gray},
      parent anchor=east,
      child anchor=west,
      l sep=14pt,
      s sep=3pt,
      anchor=west,
      inner sep=3pt,
      edge path={
        \noexpand\path[\forestoption{edge}]
        (!u.parent anchor) -- +(7pt,0) |- (.child anchor)\forestoption{edge label};
      },
    },
  },
}

% Стили узлов для линейных схем и блок-схем
\tikzset{
  box/.style={
    draw, rounded corners=2pt, align=center, font=\small,
    inner sep=4pt, minimum height=8mm
  },
  boxw/.style={box, text width=#1},
  term/.style={box, rounded corners=8pt},
  note/.style={
    draw, dashed, align=left, font=\footnotesize\itshape,
    inner sep=3pt, fill=yellow!8
  },
  flowarrow/.style={-Latex, thick},
}

% Окружение для схемы: всегда отдельным блоком на всю ширину,
% никогда не рядом с другой схемой
\newenvironment{shema}[2]
  {\begin{figure}[htbp]\centering\def\shemacaption{#1}\def\shemalabel{#2}}
  {\caption{\shemacaption}\label{\shemalabel}\end{figure}}
